\section{A universal host maze}\label{sec:universal-host}
\label{subsec:universal-host-construction}

Joining traps can change a queried compass mask even when the new edge was never
traversed. A universal induced maze therefore requires protection of the entire
execution, not just its recurrent edges.

A five-cell path already shows the issue.  Number a vertical path $0,1,2,3,4$ from
bottom to top and start the agents at $0$ and $4$.  There are legal completions with
straight-axis rows for which the complete spatial words are
\[
  (0,1,2,1)^\infty,
  \qquad
  (4,3)^\infty.
\]
The distances are $4,2,2,2$ and then repeat, but every old cell is visited.  Any proper
connected induced extension changes the mask of some visited old cell.  Thus one
cannot demand that each original trap already contain a clean attachment site.

\subsection{Clean attachment}
\label{subsubsec:clean-attachment}

\begin{lemma}[Clean attachment]
\label{lem:clean-attachment}
Let $V$ be an induced grid maze and let
\[
  (x_t,q_t,y_t,r_t)_{t\ge0}
\]
be a complete trapped execution in $V$.  Put
\[
  R=\{x_t:t\ge0\}\cup\{y_t:t\ge0\}.
\]
Let $C$ be a finite set of new cells, disjoint from $V$, such that the induced extension
$V'=V\cup C$ is a maze and
\[
  \operatorname{dist}_1(C,R)\ge2.
\]
Then the tagged execution from the same starts is identical in $V$ and $V'$.
In particular the extension introduces no rendezvous.
\end{lemma}

\begin{proof}
Induct on time.  Suppose that both tagged positions agree through time $t$.  Each
current cell belongs to $R$, and no new cell is a grid neighbor of it.  Its complete
compass mask is therefore unchanged.  The controller queries the same row, returns the
same direction and next state, and the selected old neighbor is still present.  Hence
both next tagged positions agree.  The conclusion follows simultaneously for the two
agents.
\end{proof}

The set $R$ must include transient visits. New edges at unvisited old cells are
harmless: the lemma protects queried masks, not every old mask.

\begin{figure}[t]
\centering
% Schematic only. R and its grid distance-one halo are generated explicitly.
\begin{tikzpicture}[maze,x=6.9mm,y=6.9mm]
 \path[protected] (-1.46,-0.46) rectangle (-0.54,0.46);
 \path[protected] (-0.46,-1.46) rectangle (0.46,-0.54);
 \path[protected] (-0.46,-0.46) rectangle (0.46,0.46);
 \path[protected] (-0.46,0.54) rectangle (0.46,1.46);
 \path[protected] (0.54,-1.46) rectangle (1.46,-0.54);
 \path[protected] (0.54,-0.46) rectangle (1.46,0.46);
 \path[protected] (0.54,0.54) rectangle (1.46,1.46);
 \path[protected] (1.54,-1.46) rectangle (2.46,-0.54);
 \path[protected] (1.54,-0.46) rectangle (2.46,0.46);
 \path[protected] (1.54,0.54) rectangle (2.46,1.46);
 \path[protected] (1.54,1.54) rectangle (2.46,2.46);
 \path[protected] (2.54,-0.46) rectangle (3.46,0.46);
 \path[protected] (2.54,0.54) rectangle (3.46,1.46);
 \path[protected] (2.54,1.54) rectangle (3.46,2.46);
 \path[protected] (3.54,0.54) rectangle (4.46,1.46);
 \draw[edge,line width=1.05pt] (0,0)--(1,0)--(2,0)--(2,1)--(3,1);
 \foreach \x/\y in {0/0,1/0,2/0,2/1,3/1}{\node[cell,fill=rvInk] at (\x,\y) {};}
 \draw[quiet] (3,1)--(4,1);
 \node[gate,label=above:{$g$}] at (4,1) {};
 \draw[rvInk,line width=.9pt] (4,1)--(5,1)--(6,1)--(6,2)--(7,2);
 \foreach \x/\y in {5/1,6/1,6/2,7/2}{\node[cell] at (\x,\y) {};}
 \draw[<->,rvMid,line width=.5pt] (3,.05)--(5,.05) node[midway,below,font=\footnotesize] {$2$};
 \node[paneltitle] at (-1.0,3.0) {Protect every queried mask};
 \node[note] at (8.2,2.0) {Black cells: complete visited set $R$};
 \node[note] at (8.2,1.1) {Gray halo: $\operatorname{dist}_1(\,\cdot\,,R)\le1$};
 \node[note] at (8.2,.2) {$g$: unvisited old gate};
 \node[note] at (8.2,-.7) {New connector: $\operatorname{dist}_1(C,R)\ge2$};
 \mazeCompass{16.0}{2.1}
\end{tikzpicture}

\caption{Clean attachment: the gray distance-one halo protects every cell in the complete
visited set $R$. The new connector meets only an unvisited old gate and stays outside
the halo; the drawing illustrates the clearance condition, not a particular trap.}
\label{fig:clean-attachment}
\end{figure}

\subsection{Unused-cut stretching}
\label{subsubsec:unused-cut-stretching}

The five-cell example above still has an unused edge cut.  Stretching such a cut makes
room for a protected port without modifying the masks of the old cells.

\begin{lemma}[Unused-cut stretching]
\label{lem:unused-cut-stretching}
Let the vertical cut $x=k+\tfrac12$ cross at least one edge of an induced grid maze
$V$, and suppose that neither complete trapped trajectory ever traverses an edge that
crosses this cut.  Define
\[
\phi(x,y)=
\begin{cases}
(x,y),&x\le k,\\
(x+4,y),&x\ge k+1.
\end{cases}
\]
For every old crossing edge $(k,y)(k+1,y)$ insert the four cells
\[
 (k+1,y),(k+2,y),(k+3,y),(k+4,y)
\]
between its image endpoints $(k,y)$ and $(k+5,y)$.  Then every old image cell has
exactly its original compass mask, every actually selected move of the trapped
execution joins the images of the same two old cells, and Manhattan distances between
old image cells do not decrease.  Moreover every inserted cell in the middle column
$x=k+2$ has Manhattan distance at least two from every old image cell.
\end{lemma}

\begin{proof}
Within either half-plane, $\phi$ is a translation, so noncrossing adjacencies and
compass directions are unchanged.  At a crossing edge the left endpoint still has an
east neighbor and the right endpoint still has a west neighbor, now through the
inserted straight corridor; no inserted cell is adjacent to any other old image.
Thus every old mask is exactly preserved.

Mask preservation would not by itself preserve an execution that actually used a
crossing edge, because such a move would enter a new corridor cell.  The unused-cut
hypothesis excludes this: every selected old edge lies wholly in one half-plane and
therefore maps to the same old adjacency.  Induction on time transports the entire
tagged execution.

Distances on one side are unchanged.  For points on opposite sides the horizontal
separation increases by four, so Manhattan distance cannot decrease.  Finally, old
left images have $x\le k$ and old right images have $x\ge k+5$; a cell with $x=k+2$
is therefore at horizontal distance at least $2$ from the former and at least $3$
from the latter.  This is the protected lane required for Lemma~\ref{lem:clean-attachment}.
\end{proof}

The horizontal version is obtained by writing the same geometric construction with
$x$ and $y$ interchanged.  This is a statement about the embedding of a fixed execution;
it does \emph{not} rotate or reflect the fixed controller.

For the cuts used below exactly one old tree edge crosses the cut.  The replacement
therefore subdivides that edge, and a simple route from the protected middle lane to an
external port preserves the tree property and maximum degree three.

\begin{figure}[t]
\centering
% Exact local grid scale, with the cut and its image distinguished.
\begin{tikzpicture}[maze,x=8mm,y=8mm]
 \begin{scope}
  \node[paneltitle] at (-1,2.55) {BEFORE};
  \draw[edge] (-1,0)--(0,0)--(1,0)--(2,0);
  \foreach \x in {-1,0,1,2}{\node[cell] at (\x,0) {};}
  \draw[rvMid,densely dashed] (.5,-.65)--(.5,1.05);
  \node[font=\scriptsize,above] at (.5,1.05) {$x=k+\tfrac12$};
  \node[font=\footnotesize,below] at (-.1,-.12) {$k$};
  \node[font=\footnotesize,below] at (1.1,-.12) {$k+1$};
  \node[note] at (-1,-1.15) {Unused crossing edge};
 \end{scope}
 \draw[flow,line width=.8pt] (3.25,.4)--(4.4,.4);
 \begin{scope}[shift={(6.1,0)}]
  \node[paneltitle] at (-1,2.55) {AFTER};
  \node[note] at (1.4,2.55) {Right half translated by $4$};
  \fill[rvGray!85] (1.65,-.43) rectangle (2.35,1.5);
  \draw[edge] (-1,0)--(0,0)--(5,0)--(6,0);
  \foreach \x in {-1,0,5,6}{\node[cell] at (\x,0) {};}
  \foreach \x in {1,2,3,4}{\node[gate] at (\x,0) {};}
  \draw[quiet] (2,0)--(2,1.5); \node[gate] at (2,1) {};
  \node[font=\footnotesize,below] at (0,-.12) {$k$};
  \node[font=\footnotesize,below] at (2,-.12) {$k+2$};
  \node[font=\footnotesize,below] at (5,-.12) {$k+5$};
  \node[note] at (-1,-1.15) {Old masks unchanged; lane $k+2$ protected.};
 \end{scope}
 \mazeCompass{13.3}{.75}
\end{tikzpicture}

\caption{An unused crossing edge before and after width-four stretching. Four inserted
cells preserve the old endpoint masks; the lane $x=k+2$ has clearance for a protected
branch.}
\label{fig:cut-stretching}
\end{figure}

\subsection{Portability of the junction leaves}
\label{subsubsec:junction-portability}

All nonjunction leaves of the case tree are paths and admit an explicitly proved
unused cut through an edge not traversed by either complete execution.  The two final
junction leaves require separate arguments.

For G2, the W9 symbolic words confine one execution to the cells $B,J,A,U$ while the
other alternates on its separated two-cycle; the boundary cell $V$ is never visited.
Its outward ray is therefore a clean port.  For G1, one execution is the fixed word
\[
 B^0,D^1,F^1,D^1,F^1,\ldots
\]
and neither visits $V$ nor traverses $JB$.  The other begins $A^0,J^1$.  The G1 proof
implies the dichotomy
\[
  V\text{ is never visited}
  \qquad\text{or}\qquad
  JB\text{ is never traversed by either execution}.
\]
Indeed, a first visit to $B$ can occur only as $B^1$ and immediately repeats the
already visited tagged state $J^1$, while a visit to $V$ enters the permanent
$U^1,V^1$ two-cycle.  The same execution therefore cannot visit both $B$ and $V$.
In the first alternative we use the clean outward ray of $V$; in the second we stretch
the unused cut through $JB$.  Both G1 variants are included in the fixed host in
advance.  The choice of which one supplies the witness may depend on the controller.

\subsection{Assembly of the universal tree}
\label{subsubsec:universal-tree-assembly}

Applying these port constructions to the inverse compass transports required by the
case tree gives $172$ ported fragments: $156$ unused-cut fragments and $16$
clean-boundary fragments.  The number exceeds $144$ because one oriented test can
have more than one relevant proof provenance or attachment method, and G1 requires
both of its portable alternatives in advance.

Place the $172$ fragment boxes from left to right with two empty columns between
successive boxes.  Extend each exposed port by a vertical stalk to one top horizontal
rail, and join the stalks along that rail.  The separation prevents unintended induced
adjacencies between fragments or stalks, while each fragment has exactly one connection
to the global rail.  The local pieces are trees of maximum degree three, the stalks are
paths, and the rail joins them without cycles.  The induced union is therefore a
subcubic tree.

\begin{figure}[t]
\centering
% A schematic assembly, not the complete 5212-cell coordinate drawing.
\begin{tikzpicture}[maze,x=7.5mm,y=6.5mm]
 \draw[rvInk,line width=1.15pt] (.3,3.75)--(12.5,3.75);
 \node[paneltitle] at (.25,4.48) {One host, fixed before the controller};
 \foreach \x in {1.3,4.3,7.3,10.3}{
  \draw[lightpanel] (\x-.83,-.1) rectangle (\x+.83,1.65);
  \draw[rvInk,line width=.95pt] (\x,1.45)--(\x,3.75);
  \node[cell] at (\x,3.75) {};
  \node[gate] at (\x,1.45) {};
 }
 % Representative schematic trees inside three fragment boxes.
 \draw[edge] (.8,.4)--(1.3,.4)--(1.3,1.45); \draw[edge] (1.3,.9)--(1.85,.9);
 \draw[edge] (3.8,.4)--(4.3,.4)--(4.8,.4)--(4.8,1)--(4.3,1)--(4.3,1.45);
 \draw[edge] (9.8,.4)--(10.3,.4)--(10.3,1.45); \draw[edge] (10.3,.9)--(10.85,.9);
 \foreach \x/\y in {.8/.4,1.3/.4,1.3/.9,1.85/.9,3.8/.4,4.3/.4,4.8/.4,4.8/1,4.3/1,9.8/.4,10.3/.4,10.3/.9,10.85/.9}{\node[cell] at (\x,\y) {};}
 \node at (7.3,.75) {$\cdots$};
 \draw[rvMid,line width=.4pt] (12.5,3.75)--(12.9,3.75);
 \draw[rvMid,line width=.4pt] (10.4,2.65)--(12.9,2.65);
 \draw[rvMid,line width=.4pt] (11.15,1.4)--(12.9,1.4);
 \node[note] at (13.1,3.75) {$1{,}430$ rail cells};
 \node[note] at (13.1,2.65) {$990$ stalk cells};
 \node[note] at (13.1,1.4) {$2{,}792$ fragment cells};
 \node[note] at (.25,-.65) {$172$ fragments; separated boxes; one stalk per fragment};
 \node[font=\small\bfseries] at (8.25,-1.65) {$2792+990+1430=5212$};
 \mazeCompass{17.25}{-.25}
\end{tikzpicture}

\caption{Universal-host assembly, shown schematically. Each separated fragment joins the
rail through one stalk; the three displayed cell counts are disjoint contributions
to the exact construction.}
\label{fig:global-rail}
\end{figure}

The exact audited realization contains $2{,}792$ cells in the disjoint ported
fragments, $990$ stalk cells strictly between the fragment ports and the rail, and
$1{,}430$ rail cells.  Hence
\begin{equation}
\label{eq:universal-maze-count}
  |U|=2792+990+1430=5212.
\end{equation}
Its induced graph has $5{,}211$ edges and maximum degree $3$.

\begin{theorem}[One universal two-state maze]
\label{thm:universal-5212}
There exists a fixed finite connected induced grid tree $U$ with $5{,}212$ cells and
maximum degree three such that for every legal deterministic controller $A$ with at
most two states there are state-$0$ starts $x_A,y_A\in U$, with
$\|x_A-y_A\|_1\ge2$, for which the two copies of $A$ never reach Manhattan distance
at most one.  Equivalently, the quantifier order is
\[
  \boxed{\exists U\;\forall A\;\exists(x_A,y_A).}
\]
\end{theorem}

\begin{proof}
Fix $A$.  The case tree and Theorem~\ref{thm:human-finite-basis} choose a
terminal source trap together with one simultaneous compass normalization.  Its
inverse-transported oriented instance is represented among the fixed ported fragments.
For a path leaf, its proved full-execution unused cut gives the protected embedding;
for G2 use the unvisited gate; for G1 use the clean-$V$/unused-$JB$ dichotomy.  The
local attachment lemmas preserve the complete unsuccessful execution.  The separated
global placement adds no new neighbor to any cell visited by that execution, so it is
preserved once more in the same fixed maze $U$.  Only the choice of fragment and its
recorded start pair depends on $A$.
\end{proof}

The obstruction basis and portability lemmas prove existence constructively.
The exact $5{,}212$-cell placement is separately validated from all $172$ fragment
records, their starts and embeddings, and the final induced graph; coordinate
verification does not replace the assembly proof.

\subsection{The universal-maze parameter}
\label{subsec:universal-maze-parameter}

Theorem~\ref{thm:universal-5212} suggests a parameter different from the
controller-dependent quantity $F(2)$.  Let
\[
\begin{split}
 u_2=\min\bigl\{|U|:\;& U\text{ is a finite connected induced grid maze and }\\
 &\forall A,\ |Q_A|\le2,\ \exists x_A,y_A\in U\text{ trapping }A\bigr\},
\end{split}
\]
where the starts are state-$0$ and initially noncontacting.  Theorem~\ref{thm:universal-5212}
shows that this minimum is finite and gives $u_2\le5212$.

\begin{corollary}[Every exact size above the construction]
\label{cor:universal-exact-size}
For every integer $n\ge5212$ there exists a universal induced grid tree of maximum
degree three with exactly $n$ cells.
\end{corollary}

\begin{proof}
Extend the exposed right end of the top rail by an induced horizontal ray of exactly
$n-5212$ new cells.  In the explicit placement the first new cell is at Manhattan
distance $8$ from the union of all copied original gadget-cell images, a superset of
every preserved witness trajectory; distances only increase farther along the ray.
Thus no mask queried by any preserved witness changes.  The first new cell has only the
old rail endpoint as an old neighbor, so the extension is a path attached at one
vertex.  Connectedness, acyclicity, inducedness, and maximum degree three are
preserved.
\end{proof}

The opposite direction is currently computational.

\begin{proposition}[Computer-assisted lower bound]
\label{prop:u2-lower-eight}
The universal-maze parameter satisfies
\[
  u_2\ge8.
\]
Consequently
\[
  \boxed{8\le u_2\le5212.}
\]
\end{proposition}

\paragraph{Verification basis.}
The lower bound is a quantifier-reversal statement.  It does not follow from the fact
that every individual two-state controller has some trap on at most seven cells.
Instead, an exhaustive checker enumerates every connected fixed induced grid cell set
through size seven and tests the eight compass conjugates of the extremal controller
$A^\star$.  The controller $A^\star$ has no trap below seven cells.  Among all
seven-cell geometries, a single fixed geometry traps at most four of the eight
conjugates, even when its start pair may be chosen separately for each conjugate.
Thus no maze with at most seven cells is universal for this eight-controller subset.
This proposition is computer-assisted; no deductive classification of all seven-cell
universal candidates is claimed.

It is also useful to isolate the tree-restricted version
\[
 u_{2,\mathrm{tree3}}
 =\min\bigl\{|U|:\ U\text{ is universal and is an induced grid tree with }
 \Delta(U)\le3\bigr\}.
\]
The construction proves
\[
  8\le u_2\le u_{2,\mathrm{tree3}}\le5212,
\]
while neither minimum is known.

The displayed universal-host quantifiers concern distance-one failure, not the
classical exploration-trap parameter of Section~\ref{sec:related}.
\paragraph{Sharpness data for $A^\star$.}
Minimal seven-cell traps of the extremal controller are not all trees; the exhaustive
supplementary data also contain cyclic mechanisms.  No detailed classification of
those traps is needed for Theorem~\ref{thm:universal-5212}.
